% Root and agents B and D independently passed the full literal proof.
\subsection{Extreme corners and projective endpoint coinvariants}

\begin{proposition}[Finite-target removal of an extreme corner]
\label{C:finite-target-extreme-corner}
Let $C$ be a finite category with an object $b$ having no arrows
to any other object. Put $N=C(b,b)$ and let $D$ be the full
subcategory on the other objects. Suppose the trivial right
$\mathbb QN$-module is projective. Replace $N$ by the trivial
monoid and each incoming hom-set $C(d,b)$ by its left $N$-orbit
set, leaving $D$ unchanged. This defines an actual quotient
functor $Q:C\to\bar C$.
Let $G$ be the minimal local group of $N$. Then $BQ$ is a
rational cohomology-ring equivalence and an $\mathbb F_p$
cohomology-ring equivalence whenever $p\nmid|G|$. For every
connected finite CW complex $F$, precomposition gives a weak
equivalence
\[
 \operatorname{Map}(B\bar C,F)\longrightarrow
 \operatorname{Map}(BC,F).
\]
The quotient is surjective on fundamental groups. If $BC$ has
finite CW type, $BQ$ has a homotopy left inverse.
The dual statement for an object with no incoming arrows uses
projectivity of the trivial left $\mathbb QN$-module and right
orbit sets.
\end{proposition}
\begin{proof}
For each $d\in D$ put $S(d)=C(d,b)$. Composition makes $S$ a
contravariant functor from $D$ to finite left $N$-sets. The
mixed-nerve decomposition has the actual natural homotopy-pushout
model
\[
 BC\simeq\operatorname{hocolim}
       \bigl(BD\longleftarrow E\longrightarrow BN\bigr).
\]
Here $E$ is the realization of the category bar whose degree-$k$
term is the disjoint union of $B(S(d_k)\mathbin{/\mkern-6mu/}N)$
over composable strings $d_0\to\cdots\to d_k$ in $D$.
The maps to $BD$ forget the state and its $N$-action string;
the maps to $BN$ retain that action string. This is the same
mixed-nerve model as in
Theorem~\ref{C:finite-target-opposite-corner}, with an arbitrary
category in place of the one-object category $M$.

Choose the minimal idempotent $e$ of
Lemma~\ref{C:projective-coinvariants-minimal-column}.
For every $d$, Lemma~\ref{C:opposite-corner-action-groupoid}
gives an actual functor
\[
 S(d)\mathbin{/\mkern-6mu/}N
   \longrightarrow eS(d)\mathbin{/\mkern-6mu/}G.
\]
These functors commute strictly with precomposition by every arrow
of $D$; the comma contractions have exactly the same naturality.
The groupoid has component set $N\backslash S(d)$ and finite
stabilizer groups contained in $G$. Empty incident hom-sets cause
no exception: both corresponding spaces are empty.

The maps to these component sets are rational and good-prime
homology equivalences. The category bar spectral sequence and the
pushout prove the asserted cohomology equivalences; the actual
functor makes them ring equivalences.
For a finite CW target $F$, Miller's theorem makes the canonical
constant-map inclusions from the mapping spaces of the component
sets weak equivalences onto the mapping spaces of these groupoids.
They are natural for the whole $D$-diagram. Taking the derived
homotopy limits that map the category bars into $F$, and then the
homotopy pullbacks for the pushouts, proves the mapping-space
assertion. No choice of equivariant basepoints or contractions is
needed. This repeats the finite-target proof of
Theorem~\ref{C:finite-target-opposite-corner} objectwise.

The quotient fixes the objects and is onto each hom-set. Every
edge loop in its nerve therefore lifts to an edge loop with the
same sequence of objects, proving fundamental-group surjectivity.
If $BC$ has finite CW type, surjectivity on mapping-space components
for a finite model of $BC$ produces the homotopy left inverse.
Taking opposite categories proves the dual statement.
\end{proof}

\begin{corollary}[The finite-target triangle after endpoint removal]
\label{C:finite-target-semisimple-endpoints}
Let $C$ have three linearly ordered objects, no backward arrows,
and all three forward hom-sets nonempty. Suppose the trivial left
module at the initial endpoint and the trivial right module at the
final endpoint are projective over their rational monoid algebras,
and $BC\simeq\mathbb CP^2$. Collapse both endpoints to their orbit
sets as in the preceding proposition. Write $M$ for the unchanged
middle monoid and write $s,t,u$ for the three resulting forward
hom-set cardinalities, with $u$ the long hom-set.
Then the triangle-compression monoid
\[
 P=M\sqcup(S\times T),\qquad |P|=|M|+st,
\]
has the following properties:
\begin{enumerate}
\item $BP$ is simply connected and has rational homology of
$\mathbb CP^2\vee\bigvee_{u-1}S^2$;
\item every integral degree-two class on $BP$ restricts to zero
along every map from a finite group;
\item its rational degree-two cup pairing has rank one, and it
has an integral degree-two class whose square is nonzero and
primitive modulo torsion in degree four.
\end{enumerate}
No finite-CW-type assertion about $BP$ is included.
\end{corollary}
\begin{proof}
Apply the dual extreme-corner proposition at the initial object
and the extreme-corner proposition at the final object. The orbit
sets and all their remaining actions commute with the other
endpoint quotient, so both operations are actual functors.
Let $Y$ be the classifying space of the resulting category.
It is simply connected by fundamental-group surjectivity, has
the rational cohomology ring of $\mathbb CP^2$, and the composite
quotient has a homotopy left inverse $h:Y\to BC$ by the
mapping-space assertion. As in
Corollary~\ref{C:finite-target-corner-CP2}, this gives an integral
generator $v$ with nonzero primitive free square and nullity of
every integral degree-two group restriction.

The long-hom collapse of
Proposition~\ref{C:linear-long-hom-collapse} is valid for any
simply connected $Y$, without finiteness. It gives an actual
equivalence under $Y$
\[
 BP\simeq Y\vee\bigvee_{u-1}S^2.
\]
Compose an inverse equivalence with
$h\vee1$ and a chosen finite model
$BC\simeq\mathbb CP^2$. The resulting map to
$\mathbb CP^2\vee\bigvee_{u-1}S^2$ is an isomorphism on
integral degree-two cohomology. Miller's theorem therefore gives
the asserted nullity for every such class. The rational pairing
and the primitive free square follow from the actual wedge
decomposition and the corresponding assertions for $Y$.
\end{proof}

\begin{theorem}[Projective endpoint coinvariants through sixteen arrows]
\label{C:semisimple-endpoints-sixteen}
Let $C$ have three linearly ordered objects, no backward arrows,
and nonempty forward hom-sets. Suppose the trivial left module at
the initial endpoint and the trivial right module at the final
endpoint are projective over their rational monoid algebras.
If $C$ has at most sixteen arrows, then
$BC\not\simeq\mathbb CP^2$. In particular this applies when both
endpoint algebras are semisimple over $\mathbb Q$. There is no
separate bound on the middle monoid or the individual hom-set sizes.
\end{theorem}
\begin{proof}
Suppose otherwise and apply the endpoint-removal corollary. If the
middle monoid has at most four elements, apply
Corollary~\ref{C:semisimple-ends-middle-four} to the collapsed
category, whose endpoint algebras are now $\mathbb Q$. This
contradicts its bounded rational homology and nonzero fourth
homology. Thus write $m=|M|\geq5$.
Orbit sets can only decrease cardinalities, so
\[
 m+s+t+u+2\leq16,\qquad u\geq1.
\]
The triangle-compression monoid has minimal band $s\times t$.
If either $s$ or $t$ is at most two, its rational degree-two
products vanish by the balanced presentation and
Corollary~\ref{cor:B-two-row-general}. Their pullbacks to $Y$
would vanish, contradicting its projective-plane square.
Thus $s,t\geq3$.

The numerical inequality now leaves exactly the following
possibilities, up to opposition:
\[
\begin{array}{c|c|c}
 (s,t)&\text{middle and long-arrow bounds}&\text{obstruction}\\\hline
 (3,3)&m\geq5,\ u\leq8-m\leq3&
            \text{all-rank three-by-three zero products}\\
 (3,4)&m\geq5,\ u\leq7-m\leq2&
            \text{rank-two zero products or rank-one parity}\\
 (3,5)&m=5,\ u=1&\text{three-by-five source of order }20\\
 (4,4)&m=5,\ u=1&\text{four-by-four source of order }21.
\end{array}
\]
For the first two rows, apply
Lemma~\ref{C:small-band-finite-target-cups} to the actual
compression. Its rational second-cohomology rank is $u$ and
all integral group restrictions vanish by the preceding corollary.
The rank-one case has a nonzero primitive square, so its even-square
exception is excluded as well. For the last two rows, use
Propositions~\ref{B:three-by-five-finite-target-transfer} and
\ref{B:four-by-four-finite-target-transfer}, respectively.
Here $u=1$, so the entire rational homology is that of
$\mathbb CP^2$, including vanishing third homology. All stated
transfer hypotheses are therefore supplied by the actual maps,
not by an inferred finite type of the compression.
Each possibility is impossible, proving the theorem.
\end{proof}
